% -*- coding: utf-8 -*-
% !TEX program = xelatex
\documentclass[plain,noanswer,medmath]{../../template/jnuexam}
\usepackage{../../template/kysx}

\begin{document}


\examtitle{1996}{1}


\makepart{填空题}{本题共5小题，每小题3分，满分15分}


\begin{problem}
设$\lim_{x \to \infty} \left(\frac{x + 2a}{x - a} \right)^x = 8$，则$a =$ \fillin{}．
\end{problem}


\begin{problem}
设一平面经过原点及点$(6, - 3,2)$，且与平面$4x - y + 2z = 8$垂直，则此平面方程为 \fillin{}．
\end{problem}


\begin{problem}
微分方程$y'' - 2y' + 2y = \e^x$的通解为 \fillin{}．
\end{problem}


\begin{problem}
函数$u=\ln(x+\sqrt{y^2+z^2})$在$A(1,0,1)$点处沿$A$点指向$B(3,-2,2)$点方向的方向导数为 \fillin{}．
\end{problem}


\begin{problem}
设$A$是$4 \times 3$矩阵，且$A$的秩$r(A) = 2$，而$B = \begin{pmatrix}
   1 & 0 & 2 \\
   0 & 2 & 0 \\
  -1 & 0 & 3
\end{pmatrix}$，则$r(AB) =$ \fillin{}．
\end{problem}


\makepart{选择题}{本题共5小题，每小题3分，满分15分}


\begin{problem}
已知$\frac{(x + ay)\dx + y\dy}{(x + y)^2}$为某函数的全微分，则$a$等于\pickout{}
\begin{abcd}
\item -1．        
\item 0．        
\item 1．        
\item 2．
\end{abcd}
\end{problem}


\begin{problem}
设$f(x)$有二阶连续导数，且$f'(0) = 0$,$\lim_{x \to 0} = \frac{f''(x)}{|x|} = 1$，则\pickout{}
\begin{abcd}
\item $f(0)$是$f(x)$的极大值．
\item $f(0)$是$f(x)$的极小值．
\item $(0,f(0))$是曲线$y = f(x)$的拐点．
\item $f(0)$不是$f(x)$的极值，$(0,f(0))$也不是曲线$y = f(x)$的拐点． 
\end{abcd}
\end{problem}


\begin{problem}
设$a_n > 0$（$n = 1,2, \cdots$），且$\sum_{n = 1}^{\infty}a_n$收敛，
常数$\lambda \in \left(0,\frac{\pi}{2} \right)$，
则级数$\sum_{n = 1}^{\infty}(- 1)^n\left(n\tan \frac{\lambda}{n} \right) a_{2n}$\pickout{}
\begin{abcd}
\item 绝对收敛．   
\item 条件收敛．   
\item 发散．   
\item 收敛性与$\lambda$有关．
\end{abcd}
\end{problem}


\begin{problem}
设$f(x)$有连续的导数，$f(0) = 0$,$f'(0) \ne 0$,$F(x) = \int_0^x (x^2 - t^2)f(t)\dt$，且当$x \to 0$时，$F'(x)$与$x^k$是同阶无穷小，则$k$等于\pickout{}
\begin{abcd}
\item $1$．          
\item $2$．          
\item $3$．          
\item $4$．
\end{abcd}
\end{problem}


\begin{problem}
四阶行列式$\begin{vmatrix}
  a_1 &   0 &   0 & b_1 \\
    0 & a_2 & b_2 & 0 \\
    0 & b_3 & a_3 & 0 \\
  b_4 &   0 &   0 & a_4
\end{vmatrix}$的值等于\pickout{}
\begin{abcd}
\item $a_1a_2a_3a_4 - b_1b_2b_3b_4$．					
\item $a_1a_2a_3a_4 + b_1b_2b_3b_4$．
\item $(a_1a_2 - b_1b_2)(a_3a_4 - b_3b_4)$．			
\item $(a_2a_3 - b_2b_3)(a_1a_4 - b_1b_4)$．
\end{abcd}
\end{problem}


\makepart{解答题}{本题共2小题，每小题5分，满分10分}


\begin{problem}
求心形线$r = a(1 + \cos \theta)$的全长，其中$a > 0$是常数．
\end{problem}


\begin{problem}
设$x_1 = 10$，$x_{n + 1} = \sqrt{6 + x_n}$（$n = 1,2, \cdots$），
试证数列$\left\{x_n \right\}$极限存在，并求此极限．
\end{problem}


\makepart{计算题}{本题共2小题，每小题6分，满分12分}


\begin{problem}
计算曲面积分$\iint_S (2x + z)\dy\dz + z\dx\dy,$
其中$S$为有向曲面$z = x^2 + y^2$（$0 \le z \le 1$），其法向量与$z$轴正向的夹角为锐角．
\end{problem}


\begin{problem}
设变换$\begin{cases}
 u = x - 2y, \\
 v = x + ay \\
\end{cases}$可把方程
$$6\frac{\pd^2z}{\pd x^2} + \frac{\pd^2z}{\pdx\pd y} - \frac{\pd^2z}{\pd y^2} = 0$$
化简为$\frac{\pd^2 z}{\pdu\pd v} = 0$,
求常数$a$，其中$z = z(x,y)$有二阶连续的偏导数．
\end{problem}


\makepart{}{本题满分7分}

\begin{problem}
求级数$\sum_{n = 2}^{\infty}\frac{1}{(n^2 - 1)2^n}$的和．
\end{problem}


\makepart{}{本题满分7分}

\begin{problem}
设对任意$x > 0$，曲线$y = f(x)$上点$(x,f(x))$处的切线在$y$轴上的截距等于
$\frac{1}{x}\int_0^x f(t) \dt$，求$f(x)$的一般表达式．
\end{problem}


\makepart{}{本题满分8分}

\begin{problem}
设$f(x)$在$[0,1]$上具有二阶导数，且满足条件$|f(x)| \le a$,$|f''(x)| \le b$，其中$a,b$都是非负常数，$c$是$(0,1)$内任一点，证明$|f'(c)| \le 2a + \frac{b}{2}$．
\end{problem}


\makepart{}{本题满分6分}

\begin{problem}
设$A = E - \xi \xi^T$，其中$E$是$n$阶单位矩阵，$\xi$是$n$维非零列向量，$\xi^T$是$\xi$的转置，证明：\par
\begin{enumerate*}
\item $A^2 = A$的充要条件是$\xi^T\xi = 1$；
\item 当$\xi^T\xi = 1$时，$A$是不可逆矩阵．
\end{enumerate*}
\end{problem}


\makepart{}{本题满分8分}

\begin{problem}
已知二次型$f(x_1,x_2,x_3) = 5x_1^2+ 5x_2^2+ cx_3^2- 2x_1x_2 + 6x_1x_3 - 6x_2x_3$的秩为$2$．
\begin{enumerate}
\item 求参数$c$及此二次型对应矩阵的特征值；
\item 指出方程$f(x_1,x_2,x_3) = 1$表示何种二次曲面．
\end{enumerate}
\end{problem}


\makepart{填空题}{本题共2小题，每小题3分，满分6分}


\begin{problem}
设工厂$A$和工厂$B$的产品的次品率分别为$1\%$和$2\%$，
现从由$A$和$B$的产品分别占$60\%$和$40\%$的一批产品中随机抽取一件，发现是次品，则该次品属$A$生产的概率是 \fillin{}．
\end{problem}


\begin{problem}
设$\xi$, $\eta$是两个相互独立且均服从正态分布$N\left(0,\left(\tfrac{1}{\sqrt{2}}\right)^2\right)$的随机变量，
则随机变量$|\xi-\eta|$的数学期望$E(|\xi-\eta|) =$ \fillin{}．
\end{problem}


\makepart{}{本题满分6分}

\begin{problem}
设$\xi$, $\eta$是相互独立且服从同一分布的两个随机变量，已知$\xi$的分布律为
$$P\left\{\xi = i \right\} = \frac{1}{3}, \quad i=1,2,3.$$
又设$X = \max(\xi ,\eta)$, $Y = \min(\xi ,\eta)$．
\begin{enumerate}
\item 写出二维随机变量$(X,Y)$的分布律：
\par\centerline{\begin{tabularx}{0.5\linewidth}{|c|*{3}{X|}}
\hline
\diagbox{$X$}{$Y$} & $1$ & $2$ & $3$ \\
\hline
$1$ & & & \\
\hline
$2$ & & & \\
\hline
$3$ & & & \\
\hline
\end{tabularx}}
\item 求随机变量$X$的数学期望$E(X)$．
\end{enumerate}
\end{problem}


\end{document}
